\documentclass[10pt,a4paper]{amsart}
\usepackage[margin=19mm]{geometry}
\usepackage{amsmath,amssymb,mathtools}
\usepackage[scr=rsfs,cal=euler]{mathalfa}
\usepackage{xfrac}
\usepackage[unicode,hidelinks,bookmarks=false]{hyperref}
\newcommand{\ie}{i.e.\ }
\newcommand{\cf}{cf.\ }
\newcommand{\resp}{respectively\ }
\newcommand{\Subsection}{\subsection}
\providecommand{\nobreakdash}{\nobreak\hbox{-}}
\setlength{\emergencystretch}{2em}
\multlinegap0pt
\usepackage{mathtools}
\usepackage{amssymb}
\usepackage[shortlabels]{enumitem}
\usepackage[normalem]{ulem}
\newtheorem{thm}{Theorem}
\newtheorem{lem}{Lemma}
\newtheorem{prop}{Proposition}
\newcommand{\Gal}{\mathrm{Gal}}
\newcommand{\Tr}{\mathrm{Tr}}
\newcommand{\Z}{\mathbb Z}
\newcommand{\notmid}{\mathrel{\ooalign{$\mkern-5mu\not$\crcr$|$}}}
\title{Multiplicative $f$-ic forms on algebraic varieties arising from Thaine's generalized Jacobi sums}
\author{Akinari Hoshi, Kazuki Kanai}
\date{Source: arXiv:2605.05039v1 (CC BY 4.0)}
\begin{document}
\maketitle

\noindent\emph{Source and licence.} Multiplicative $f$-ic forms on algebraic varieties arising from Thaine's generalized Jacobi sums. Version arXiv:2605.05039v1 (CC BY 4.0), distributed by arXiv under the Creative Commons Attribution 4.0 International licence, http://creativecommons.org/licenses/by/4.0/, as recorded in the arXiv licence metadata for this exact version and shown on the abstract page https://arxiv.org/abs/2605.05039v1. Full licence notice: \texttt{licenses/arxiv-2605.05039-CC-BY-4.0.txt}.\par\smallskip

\noindent\textit{Editorial scope.} Counted unit: the article's lem lem:h\_regularity (source0.tex lines 2432--2434). The article's complete original proof of it is delivered with it (source0.tex lines 2436--2529, one \texttt{proof} environment of 94 lines). No mathematics is written, repaired or re-derived: the statement and the proof are the article's own lines, in the article's own notation, and no retained line is substituted or re-flowed.  Retained uncounted as the source-local context the proof needs: lines 393--424; lines 2270--2302; lines 2405--2413.  Nothing was omitted from any retained span, so the package carries no omission record.  No retained line contains a citation, so the wrapper carries no bibliography and no bibliography entry.  No retained line contains a figure environment, an image-inclusion command or drawing code, so no image is delivered and the package declares no asset dependency.  Editorial formatting only, in addition: an AMS \texttt{amsart} preamble that restates the article's own declarations and macros used by the retained lines, namely the environments \texttt{thm}, \texttt{lem}, \texttt{prop} on separate counters, together with the standard packages those lines need.  The retained environments print in this document as Theorem 1, Lemma 1, Proposition 1 in order of appearance, whereas the article prints the same items with the numbering of its own full document.  The complete native source of the article as distributed in the arXiv e-print for this exact version is bundled unchanged as \texttt{sources/199917/source0.tex}.
\begin{thm}\label{mainth3}
Let $f\geq2$ be an integer and $l$ be an odd prime with $l \equiv 1\ ({\rm mod}\ f)$. 
Write $l=ef+1$.
Let $K$ be a field with char $K \notmid f!$, $l$ and 
$\Gal(K(\zeta_l)/K)\simeq (\Z/l\Z)^{\times}=\langle \gamma \rangle$. 
Define $\beta := \gamma^e$ and $M=L^{\langle \beta \rangle}$ where $L=K(\zeta_l)$, 
i.e., $\Gal(L/M)\simeq \Z/f\Z = \langle \beta \rangle \leq  (\Z/l\Z)^{\times}$.
Let $V \subset K^{l-1}$ be
an algebraic $K$-variety defined by
\begin{align*}
h_m({\bf x})=0\quad (m = 1,\ldots,e-1)
\end{align*}
where
\begin{align*}
h_m({\bf x})&=h_m(x_1,\ldots,x_{l-1})\\
&=\left(\sum_{i_0 + \beta i_1 + \cdots +\beta^{f-1}i_{f-1} \equiv \gamma^m \ (l)}x_{i_0}x_{i_1}\cdots x_{i_{f-1}}\right) - \left(\sum_{i_0 + \beta i_1 + \cdots +\beta^{f-1}i_{f-1} \equiv \gamma^{m+1} \ (l)}x_{i_0}x_{i_1}\cdots x_{i_{f-1}}\right)
\end{align*}
and each subscript of the $x_{i}$'s should be taken modulo $l$ with the convention $x_0=0$.
Let $h({\bf x})=h(x_1,\ldots,x_{l-1})$ be a regular $f$-ic form defined by
\begin{align*}
h({\bf x})
&=\left(\sum_{i_0 + \beta i_1 + \cdots +\beta^{f-1}i_{f-1}\equiv 0 \ (l)}x_{i_0}x_{i_1}\cdots x_{i_{f-1}}\right)
-\left(\sum_{i_0 + \beta i_1 + \cdots +\beta^{f-1}i_{f-1}\equiv \gamma \ (l)}x_{i_0}x_{i_1}\cdots x_{i_{f-1}}\right).
\end{align*}
Then $h({\bf x})$ is multiplicative on $V$.
Moreover, for $d \in (\Z/l\Z)^{\times}$, $\varphi_d ({\bf x},{\bf y})=(z_{1,d},\ldots,z_{l-1,d}) \in V$ 
can be given explicitly by
\[
z_{k,d}=(-1)^{f-1} \left(\sum_{i=1}^{l-1}x_{-d^{-1}i}y_{k-i} - \sum_{i=1}^{l-1}x_{-d^{-1}i}y_{-i}\right)\quad (k=1,\ldots,l-1)
\]
where each subscript of the $x_i$'s and the $y_j$'s should be taken modulo $l$ with the convention $x_0=y_0=0$.
\end{thm}

\begin{lem} \label{lem:BEW2}
Let $f\geq2$ be an integer and $l$ be an odd prime with $l \equiv 1\ ({\rm mod}\ f)$. 
Write $l=ef+1$.
Let $K$ be a field with char $K\notmid f!$, $l$ and 
$\Gal(K(\zeta_l)/K)\simeq (\Z/l\Z)^{\times}=\langle \gamma \rangle$. 
Define $\beta := \gamma^e$ and $M=K(\zeta_l)^{\langle \beta \rangle}$, 
i.e., $\Gal(K(\zeta_l)/M)\simeq \Z/f\Z = \langle \beta \rangle \leq  (\Z/l\Z)^{\times}$.
Let ${\bf p}\in K^{\times}$ and
\[
\alpha_{\bf p}=\sum_{k=0}^{l-1}a_k\zeta_l^k \quad (a_k \in K)\\
\quad \text{with} \quad
N_{K(\zeta_l)/M}(\alpha_{\bf p})= {\bf p}.
\]
Then we have
\begin{enumerate}
\item [$(1)$] ${\bf p} =S_0 - S_1$,
\item [$(2)$] $S_1=S_2=\cdots =S_{e}$
\end{enumerate}
where 
\begin{align*}
S_m = 
\left\{
\begin{array}{ll}
\displaystyle
\sum_{i_0 + \beta i_1 + \cdots +\beta^{f-1}i_{f-1}\equiv 0 \ (l)}a_{i_0}a_{i_1}\cdots a_{i_{f-1}} & (m = 0)\\
\displaystyle
\sum_{i_0 + \beta i_1 + \cdots +\beta^{f-1}i_{f-1}\equiv \gamma^{m} \ (l)}a_{i_0}a_{i_1}\cdots a_{i_{f-1}} & (m = 1,\ldots,e)
\end{array}
\right.
\end{align*}
and each subscript of the $a_k$'s should be taken modulo $l$. 
Moreover, we may assume that $a_0=0$ without loss of generality.
\end{lem}

\begin{prop}\label{prop:norm_polarization}
Let $E/F$ be a separable field extension of degree $f$, and assume that char $F \nmid f!$. 
Let $N_{E/F}\colon E\to F$ be the norm map.
Then
\[
\theta_N(u_1,\dots,u_f):=\frac{1}{f!}\,\Delta^f N_{E/F}(u_1,\dots,u_f)
\]
is a symmetric $f$-linear form on $E$.
\end{prop}

\begin{lem}\label{lem:h_regularity}
With the notation of Theorem~\ref{mainth3}, the $f$-ic form $h({\bf x})$ on $K^{l-1}$ is regular.
\end{lem}

\begin{proof}
Set $L=K(\zeta_l)$ and $M=L^{\langle\beta\rangle}$, so $[L:M]=f$.
Via the $K$-linear isomorphism
\[
\iota:K^{l-1}\xrightarrow{\ \sim\ }L,\qquad
{\bf x}=(x_1,\dots,x_{l-1})\longmapsto
\alpha_{\bf x}:=\sum_{i=1}^{l-1}x_i\zeta_l^i,
\]
we identify $K^{l-1}$ with $L$. Let $\sigma\in\Gal(L/M)$ be the generator determined by
$\sigma(\zeta_l)=\zeta_l^\beta$. Then
\[
N_{L/M}(z)=\prod_{t=0}^{f-1}\sigma^t(z)\qquad(z\in L).
\]

\medskip\noindent
\textbf{Step 1: $h$ as a $K$-linear projection of the norm.}
As in the proof of Lemma~\ref{lem:BEW2}, expanding $N_{L/M}(\alpha_{\bf x})$ and grouping terms by
$\langle\beta\rangle$-orbits yields
\begin{equation}\label{eq:norm_period_expansion}
N_{L/M}(\alpha_{\bf x})
= S_0({\bf x})+\sum_{m=1}^{e}S_m({\bf x})\,\eta(m-1)
\end{equation}
where $S_m({\bf x})$ are the polynomials defined in Lemma~\ref{lem:BEW2} (with $x_0=0$ and indices modulo $l$).
Since $\eta(0)+\cdots+\eta(e-1)=\zeta_l+\cdots+\zeta_l^{l-1}=-1$, we can rewrite \eqref{eq:norm_period_expansion} as
\[
N_{L/M}(\alpha_{\bf x})
=\sum_{r=0}^{e-1}\bigl(S_{r+1}({\bf x})-S_0({\bf x})\bigr)\,\eta(r).
\]
Define a $K$-linear projection $\lambda:M\to K$ by
\[
\lambda(\eta(0))=-1,\qquad \lambda(\eta(r))=0\ \ (1\le r\le e-1).
\]
Then applying $\lambda$ gives
\[
\lambda\bigl(N_{L/M}(\alpha_{\bf x})\bigr)=-(S_1({\bf x})-S_0({\bf x}))=h({\bf x}),
\]
hence
\begin{equation}\label{eq:h_as_proj_norm}
h=\lambda\circ N_{L/M}\circ\iota.
\end{equation}

\medskip\noindent
\textbf{Step 2: Polarization and a trace specialization.}
Let $\theta_N$ be the symmetric $f$-linear form associated with $N_{L/M}$ over the base field $M$
(Proposition~\ref{prop:norm_polarization} applied to $E=L$, $F=M$), and let $\theta_h$ be the symmetric
$f$-linear form associated with $h$.
By \eqref{eq:h_as_proj_norm} and linearity of polarization,
\begin{equation}\label{eq:theta_h_comp}
\theta_h=\lambda\circ\theta_N\circ(\iota\times\cdots\times\iota).
\end{equation}
A standard computation from $N_{L/M}(z)=\prod_{t=0}^{f-1}\sigma^t(z)$ gives, for $\alpha,v\in L$,
\begin{equation}\label{eq:thetaN_trace}
\theta_N(v,\alpha,\dots,\alpha)
=\frac{1}{f}\sum_{k=0}^{f-1}\sigma^k(v)\prod_{j\ne k}\sigma^j(\alpha)
=\frac{1}{f}\Tr_{L/M}\!\bigl(vP(\alpha)\bigr),
\qquad
P(\alpha):=\prod_{j=1}^{f-1}\sigma^j(\alpha).
\end{equation}

\medskip\noindent
\textbf{Step 3: Regularity.}
It suffices to show that if $\alpha \in L$ satisfies $\theta_h(v,\alpha,\dots,\alpha)=0$ for all $v\in L$,
then $\alpha=0$.
Indeed, assuming this, let $u\in L$ be such that
$\theta_h(u,w_1,\dots,w_f)=0$ for all $w_1,\dots,w_f\in L$.
Then, in particular, by taking $w_1=v$ and $w_2=\cdots=w_f=u$, we obtain
$\theta_h(u,v,u,\dots,u)=0$ for all $v\in L$.
By symmetry of $\theta_h$, this implies
$\theta_h(v,u,\dots,u)=0$ for all $v\in L$.
Hence $u=0$.

Assume $\alpha\in L$ satisfies $\theta_h(v,\alpha,\dots,\alpha)=0$ for all $v\in L$.
Using \eqref{eq:theta_h_comp} and \eqref{eq:thetaN_trace} we obtain
\[
\lambda\!\left(\Tr_{L/M}(vP(\alpha))\right)=0\qquad(\forall v\in L).
\]
Set $\Psi:=\lambda\circ\Tr_{L/M}:L\to K$. 
Then $\Psi\neq 0$: indeed, $\eta(0)+\cdots+\eta(e-1)=-1$ implies $\lambda(1)=1$, hence
\[
\Psi(1)=\lambda\!\left(\Tr_{L/M}(1)\right)=\lambda(f\cdot 1)=f\neq 0
\]
since char $K\nmid f$.
Thus $\Psi(vP(\alpha))=0$ for all $v\in L$.

Because $L/K$ is finite separable, the trace pairing $(x,y)\mapsto\Tr_{L/K}(xy)$ is non-degenerate.
Therefore there exists $c\in L^\times$ such that $\Psi(z)=\Tr_{L/K}(cz)$ for all $z\in L$. 
Consequently,
\[
0=\Psi(vP(\alpha))=\Tr_{L/K}(c\,v\,P(\alpha))\qquad(\forall v\in L),
\]
and non-degeneracy forces $cP(\alpha)=0$, hence $P(\alpha)=0$. Since $L$ is a field and
$P(\alpha)=\prod_{j=1}^{f-1}\sigma^j(\alpha)$, we conclude $\alpha=0$.
Thus $\theta_h(v,\alpha,\dots,\alpha)\equiv 0$ implies $\alpha=0$, i.e.\ $h$ is regular.
\end{proof}

\end{document}
